\begin{figure}[htbp]
\centering
\begin{tikzpicture}[x=.87cm,y=.87cm,font=\small]
  \fill[paper] (0,0) rectangle (12.4,6.45);
  \draw[source!75,line width=1.1pt] (.72,2.0) .. controls (4.4,1.3)
    and (8.7,4.2) .. (11.85,3.45);
  \node[text=source,below,font=\scriptsize] at (7.8,2.15) {长江方向（示意）};
  \fill[dynasty!45] (6.1,4.86) ellipse (1.35 and .82);
  \node[align=center,text=white] at (6.1,4.86) {西晋\\洛阳中枢};
  \fill[timeline!31] (1.7,2.2) ellipse (1.2 and .75);
  \node[align=center] at (1.7,2.2) {益州\\王濬水军};
  \fill[dynasty!24] (4.45,3.2) ellipse (1.15 and .7);
  \node[align=center,font=\scriptsize] at (4.45,3.2) {荆州、襄阳\\杜预部署};
  \fill[dynasty!21] (8.3,4.0) ellipse (1.1 and .68);
  \node[align=center,font=\scriptsize] at (8.3,4.0) {淮南\\王浑方向};
  \fill[source!30] (10.6,2.35) ellipse (1.25 and .74);
  \node[align=center] at (10.6,2.35) {建业\\孙吴中枢};
  \draw[->,very thick,color=timeline] (2.75,2.35) .. controls (4.25,2.45)
    and (6.5,2.35) .. (9.3,2.38);
  \node[text=timeline,above,font=\scriptsize] at (6.0,2.68) {益州水军沿江东下};
  \draw[->,thick,color=dynasty] (5.0,3.6) .. controls (6.5,3.15)
    and (7.8,2.85) .. (9.45,2.58);
  \node[text=dynasty,above,font=\scriptsize] at (7.2,3.58) {荆州方向并进};
  \draw[->,thick,dashed,color=dynasty] (7.0,4.44) .. controls (8.23,4.05)
    and (8.93,3.55) .. (9.55,2.96);
  \node[text=dynasty,right,font=\scriptsize] at (8.43,4.47) {淮南牵制};
  \node[draw=source!75,fill=paper,rounded corners=2pt,align=center,
    inner sep=3pt,font=\scriptsize] at (6.45,.58)
    {279--280 年的合流是多处资源被组织为同一计划，不等于地方差异消失。};
\end{tikzpicture}
\caption{西晋伐吴的多路推进（279--280 年，示意）。据《晋书·武帝纪》、
王濬、杜预等列传与《资治通鉴》卷八十概括。箭线提示益州、荆州、淮南与
建业的关系，不表示舰队规模、行军日程、完整路线或精确战线。}
\label{fig:vol02b:jin-conquest-wu}
\end{figure}
